% Root and agents B and D have passed the full literal proof; independent finite gates cited below are complete.
\subsection{The exceptional six-by-four opposite-corner profile}

\begin{lemma}[The remaining six-element source is unique]
\label{C:six-local-source-normal-form}
Suppose $|M|=6$, $H^2(BM;\mathbb Q)\ne0$, and its row
augmentation is not projective over the contracted
minimal-band algebra.  Then, up to isomorphism, $M$ is
\[
 M=\{1\}\sqcup\{x_{ij}:i,j=0,1\}\sqcup\{h\},
\]
with the usual rectangular-band multiplication and
\[
 hx_{ij}=x_{0j},\qquad x_{ij}h=x_{i0},\qquad h^2=x_{00}.
\]
In particular the contracted algebra is the dual-number
algebra and the row module is its residue module.
\end{lemma}
\begin{proof}
The minimal ideal must have at least two rows and columns
and is therefore the $2\times2$ band.  The contracted
algebra has basis $1,\bar h$ with $\bar h^2$ equal to
$0,\bar h$, or $1$.  In the latter two cases it is
semisimple and the row module is projective.
In the first case the row character of $h$ is zero,
so its row map is constant; the nonzero balanced
degree-two form forces the same zero column character,
so its column map is constant as well.
Relabel the row and column constants to be zero.
All mixed products are then the displayed ones.
The profile of $h^2$ forces its band value to be
$x_{00}$, proving uniqueness.
\end{proof}

\begin{lemma}[The exceptional corner has only one remaining core pattern]
\label{C:six-four-seven-core-pattern}
Let $M$ be the preceding source and let
$N=\{1,p,q,x\}$ be the exceptional nilpotent four-element
opposite corner.  If a seven-state $(N,M)$-collage has
type $\mathbb CP^2$, its minimal-band state core consists
of exactly two full two-state orbits.  On these orbits
the three nonidentity left maps $P,Q,E$ all induce the
same graph: one vertex orbit and one loop-edge orbit,
with each sending the loop orbit to the vertex orbit.
\end{lemma}
\begin{proof}
Use the actual second-cohomology comparison of
Proposition~\ref{C:small-opposite-natural-graph-comparison}.
For the nilpotent corner its normalized graph is obtained
from the $P,Q$ graph by deleting precisely those loop
edges indexed by noncore states in $ES$.
The comparison is only a second-cohomology isomorphism;
no boundedness of its graph cofiber is inferred.

The projective-plane generator pulls back from a
rank-one nonzero square on this graph cofiber.
If its invariant edge functional were zero, the class
would pull back from an ordinary component-act cone.
Lemma~\ref{C:nine-unbounded-act-cups} would kill its square.
Thus that edge functional is nonzero.  For a band
idempotent $e$, it is determined by its restriction
to $We=H_1(\Gamma e;\mathbb Q)$, which is nonzero.

A singleton band-state orbit gives a band-fixed vertex
after applying a left retraction if necessary.  This
vertex is fixed by all $M$ and all $N$, so the fixed-state
argument gives zero products.  Consequently every band
orbit has two states and there are at most three.
The image graph $\Gamma e$ has at most one state for
each such orbit.  One orbit cannot support a cycle.

With three orbits there are six core states and one
possible transient.  Before the loop deletions, the
three-state core graph must have positive first homology,
so all its edges are loops and $P=Q$ on the band core.
If $P=Q$ everywhere, the normalized graph is a disjoint
union of bouquets.  Strict loop doubling fixes its
vertices and commutes with every actor.  The graph-component
argument of Lemma~\ref{C:graph-component-eigenvalues}
gives eigenvalues one and two on its positive cohomology;
the eigenvalue-one second classes pull back from the
ordinary component cone.  They have zero squares by
Lemma~\ref{C:nine-unbounded-act-cups}, while an
eigenvalue-two class has square of eigenvalue four
and therefore zero.  Its rank-one square is impossible.
Otherwise the sole transient $z$ has two distinct
core endpoints in the same band orbit, hence in
opposite columns.  It is a nonloop edge.

Every state in $ES$ satisfies $P=Q$, because
$PE=QE=P$.  Thus this nonloop edge is not deleted by
the normalized graph quotient.  In the graph-localization
complex of Theorem~\ref{C:two-column-graph-localization}
it is the sole surviving noncore edge, and its boundary
is the nonzero difference of the two column classes.
Again the complex is
$[\mathbb Q\xrightarrow{\pm1}\mathbb Q]$.
The graph cofiber has zero fourth rational homology
without any boundedness assumption, a contradiction.

There must therefore be two band orbits.  Positive
$H_1(\Gamma e)$ forces one to be a vertex orbit and the
other a loop-edge orbit.  On the orbit set $P,Q$ send
the second to the first.  The equality $E^2=P$ and
the fact that $E$ fixes the common left core force
$E$ to do the same: it cannot fix the loop orbit.
All commuting maps between full two-column orbits
preserve the column labels.  This proves the stated
literal normal form.
\end{proof}

\begin{lemma}[The complete exceptional seven-state biset control]
\label{C:six-four-seven-nilpotent-control}
In the normal form of the preceding lemma there are
$244$ labeled-profile right actions and $4340$ raw
left actions.  Exact state conjugacies leave $2237$
literal categories; every one either fails the required
second mod-two Betti number or has a zero square.
\end{lemma}
\begin{proof}
Label the two core orbits by $(0,1)$ and $(2,3)$.
The first is the vertex orbit, and the second is the
loop orbit.  Each of the three transient states has
a row-orbit profile in $\{0,1\}^2$; sort their multiset.
The extra right actor $h$ restricts to $(0,0,2,2)$.
For a transient $s$, its image must have constant profile
$(\gamma_s(0),\gamma_s(0))$.  Every such image is tried,
then $h^2=x_{00}$ and every source product are checked.

The left map $E$ restricts to $(0,1,0,1)$ on the
four core states.  Its transient values must be
$0,1$ or transients of profile $(0,0)$.
Try every such assignment commuting with the complete
right action and retain $E^3=E^2$.  Set $P=E^2$.
Then try every commuting idempotent $Q$ with the same
image as $P$, imposing $QE=P$.  These conditions are
exactly the left-corner equations:
\[
\begin{gathered}
 P^2=P,\quad Q^2=Q,\quad PQ=Q,\quad QP=P,\\
 E^2=P,\quad EP=P,\quad EQ=Q,\quad PE=QE=P.
\end{gathered}
\]
For example $E$ fixes the image of $P$ because
$EP=E^3=P$, so $EQ=Q$ follows for a common-image
idempotent $Q$.  All sixteen left products and all
mixed equations are also rechecked literally.
The case $P=Q$ was already excluded.

Only permutations of equal-profile transient labels
are used to identify these actions.  In particular,
$p,q$ are not interchanged: the fixed nilpotent corner
distinguishes $p=x^2$.  The producer
\path{code/a_six_nilpotent_seven_acts.py}
stores every retained right action and each $P,Q,E$,
as well as the actual mod-two category certificate.
Its complete output is
\path{code/outputs/a_six_nilpotent_seven_acts.json}.
Exactly $1198$ categories have the wrong second
Betti number; all $1039$ others have explicit
nonboundary generators and square primitives.
The standalone checker
\path{code/b_verify_six_nilpotent_seven_acts.py}
uses fresh exhaustive $7^3$ pointwise assignments for
the extra right actor and the transient left maps.
It reproduces all $244$ right actions and $4340$
raw left actions, and expands the $2237$ retained
orbits disjointly, without interchanging $p,q$.
Its independent literal bars check $1816444$ typed
associativity, $1308645$ boundary-square, $514305$
cocycle and $2711790$ square-primitive equations.
The complete output is
\path{code/outputs/b_six_nilpotent_seven_acts_independently_verified.json}.
\end{proof}

\begin{theorem}[The full seventeen-arrow six-by-four profile]
\label{C:six-four-seven-excluded}
No two-object triangular category with corner orders
six and four and seven cross arrows has type
$\mathbb CP^2$.
\end{theorem}
\begin{proof}
Use opposition to write the cross arrows from the
order-six corner to the order-four corner; the other
orientation is covered by applying the same statement
to the opposite category.
If the source has zero second rational cohomology,
both vertex restrictions vanish and the relative
suspension kills the square.
Every nonexceptional four-element opposite corner
is covered by
Corollary~\ref{C:small-opposite-six-source-regular}.
For the nilpotent exception, every projective contracted
source row is covered by
Corollary~\ref{C:small-opposite-P5-all-states}.
The remaining source is exactly
Lemma~\ref{C:six-local-source-normal-form}.
Lemmas~\ref{C:six-four-seven-core-pattern} and
\ref{C:six-four-seven-nilpotent-control} exhaust its
actual seven-state bisets and exclude them.
\end{proof}
